\documentclass[11pt]{article}
\usepackage{amsmath,amssymb,amsthm}
\usepackage{geometry}
\geometry{margin=25mm}

\title{拡張連続性公理とSolver集合、連続体仮説}
\author{}
\date{}

\newtheorem{theorem}{定理}
\newtheorem{definition}{定義}
\newtheorem{remark}{注意}

\begin{document}
\maketitle

\section{1. Solver集合と拡張連続性公理による一般連続体仮説および特異基数問題の再定式化}

\subsection{1.1 序論}

連続体仮説（Continuum Hypothesis, CH）および一般連続体仮説（Generalized Continuum Hypothesis, GCH）は、ZFC集合論の内部において長年研究されてきた中心的問題である。
特に、G\"odel および Cohen によって CH が ZFC から独立であることが示されて以降、これらの命題は「真偽未決定な命題」として扱われることが通例となった。
しかしこの扱いは、暗黙のうちに \emph{ZFC という理論を一意的な前提} と見なし、その過程を無視している。
ここではSolver集合および拡張連続性公理（Extended Continuum Axiom, ECA）の枠組みを用いることで、これが CH・GCH・特異基数問題の理解を妨げていたことを示す。

\subsection{1.2 Solver集合と公理の枠組み}
\subsubsection{1.2.1 Solver集合の定義}
Solver集合を次のように定義する：
\[
\mathcal{S} = \{ S \mid S = (\mathcal{A}_S, \mathcal{R}_S) \}
\]
ここで、$\mathcal{A}_S$ は公理の選択集合、$\mathcal{R}_S$ は推論規則および構成制約の集合である。
各 $S \in \mathcal{S}$ は公理そのものではなく、公理を定義するための集まり (集合) を表す。

\subsubsection{1.2.2 公理写像}

公理写像を次で定義する：
\[
\mathcal{T} : \mathcal{S} \to \mathrm{Th}
\]
$\mathcal{T}(S)$ は Solver $S$ によって表現される公理集合を表す。
矛盾や写像が定義できないなどの場合、$\mathcal{T}(S)$ は未定義とする。

\subsubsection{1.2.3 評価関数と拡張連続性公理}
拡張連続性公理に基づき、評価関数
\[
I : \mathcal{S} \to [0,1]
\]
を導入する。$I(S)$ は、$\mathcal{T}(S)$ が内部無矛盾性、推論閉包性、構成可能性をどの程度満たすかを測度として表す。

\subsection{1.3 ZFC の Solver 集合上での位置づけ}
ZFC を生成する Solver 全体の集合を
\[
\mathcal{Z} = \{ S \in \mathcal{S} \mid \mathcal{T}(S) \cong \mathrm{ZFC} \}
\]
と定義する。このとき重要な点は、$\mathcal{Z}$ が一般に単一要素ではなく、複数の Solver からなる集合となることである。

したがって、「ZFC」という理論は Solver 集合上の同型類として理解され、その内部命題の振る舞いは、生成に用いられた Solver 集合の公理選択に依存しうる。

\subsection{1.4 一般連続体仮説の構造的欠点}
一般連続体仮説（GCH）は次を主張する：
\[
\forall \kappa \ (\kappa \text{ が無限基数} \Rightarrow 2^\kappa = \kappa^+)
\]
これは任意の無限基数 $\kappa$ に対し、$\kappa$ とその冪集合の濃度の間に中間的基数が存在しないことを全称的に要求する命題である。
しかしこの主張は、$\kappa$ が正則基数であるか特異基数であるかを区別しない。冪集合演算 $2^\kappa$ の挙動がこの区別に強く依存するにもかかわらず、GCH は単一の規則を全ての無限基数に適用している点に構造的欠陥を持つ。

\subsection{1.5 特異基数における冪集合問題}
特異基数 $\kappa$ に対する $2^\kappa$ の濃度は、ZFC の内部では完全に決定されていない。この問題は Singular Cardinals Hypothesis (SCH) や Shelah の PCF 理論によって部分的に解明されているが、依然として一般的な決定原理は存在しない。
この未解決性は、ZFC が不完全であることを意味するのではなく、冪集合操作が \emph{極限基数に対してどのような構成規則を持つか} が理論内部で明示されていないことに起因する。

\subsection{1.6 Solver集合による統一的説明}
Solver集合の枠組みでは、冪集合演算は ZFC に固有の操作ではなく、各 Solver 集合によって定義される構成規則として扱われる。

このとき次が成立する：
\begin{itemize}
  \item 正則基数に関する冪集合の挙動は、多くの Solver において比較的一様である。
  \item 特異基数に関する冪集合の挙動は、Solver が許容する極限構成規則に強く依存する。
\end{itemize}

その結果、GCH を満たす Solver 集合と特異基数において GCH が破れる Solver 集合が同一の $\mathcal{Z}$ 内に存在するといったことが起こる。
したがってSolver集合におけるGCHの破綻や特異基数の冪集合問題は、集合の性質ではなく \emph{Solver 集合における公理選択の差異} に由来するものとして説明できる。

\subsection{1.7 結論}
本稿では、Solver集合と拡張連続性公理の枠組みを用いることで、一般連続体仮説の構造的欠点および特異基数における冪集合濃度の未解決性が、共通の原因――すなわち冪集合操作を Solver 非依存な全称命題として扱った点――に由来することを示した。
この視点において、CH や GCH はもはや単なる真偽未決定命題ではなく、Solver 集合上の構成規則を分類する命題として再解釈される。

\section{2 特異基数における冪集合写像と ZFC+GCH が Solver 集合に課す制約}

\subsection{2.1 序論}
前稿では、一般連続体仮説（GCH）および特異基数における冪集合濃度の未解決性が冪集合操作を Solver 非依存な全称命題として扱った点に起因することを示した。
本稿ではさらに一歩進め、特異基数 $\kappa$ に対する $2^\kappa$ を Solver 集合上の写像として明示的に定義し、その結果として ZFC+GCH が Solver 集合上で極端に強い制約を課すことを定理として示す。

\subsection{2.2 Solver 集合と基数生成の再確認}
Solver 集合 $\mathcal{S}$、公理写像 $\mathcal{T}$、評価関数 $I$ は既に定義されているものとする。
以降、ZFC を生成する Solver の集合を
\[
\mathcal{Z} = \{ S \in \mathcal{S} \mid \mathcal{T}(S) \cong \mathrm{ZFC} \}
\]
と表す。

\subsection{2.3 特異基数と極限構成}
\begin{definition}[特異基数]
基数 $\kappa$ が特異基数であるとは、$\kappa$ が無限基数であり
\[
\mathrm{cf}(\kappa) < \kappa
\]
を満たすことをいう。
\end{definition}

特異基数 $\kappa$ は、増大する基数列 $\langle \kappa_i \mid i < \mathrm{cf}(\kappa) \rangle$ の極限として表される。
この極限構成が冪集合演算に本質的な影響を与える。

\subsection{2.4 冪集合写像の Solver 上での定義}
\begin{definition}[Solver 上の冪集合写像]
特異基数 $\kappa$ に対し、Solver 集合上の写像
\[
\mathcal{P}_\kappa : \mathcal{Z} \to \mathrm{Card}
\]
を次で定義する：
\[
\mathcal{P}_\kappa(S) = \left| \mathcal{P}^S(\kappa) \right|
\]
ここで $\mathcal{P}^S(\kappa)$ は、Solver $S$ が許容する構成規則の下で生成される $\kappa$ の部分集合全体を表す。
\end{definition}

この定義において、$\mathcal{P}_\kappa(S)$ は集合論的対象ではなく Solver に依存する値である。

\subsection{2.5 特異基数における依存性構造}

特異基数の場合、$\mathcal{P}_\kappa(S)$ は次の要素に依存する：
\begin{itemize}
  \item $\kappa$ を表す増大基数列の選択
  \item 極限段階で許容される合成操作
  \item 下位基数 $\kappa_i$ に対する冪集合規則
\end{itemize}

この依存性は正則基数の場合には消失するが、特異基数では本質的である。

\subsection{2.6 ZFC+GCH が課す制約}
ZFC+GCH を生成する Solver の集合を
\[
\mathcal{Z}_{\mathrm{GCH}} = \{ S \in \mathcal{Z} \mid \mathcal{T}(S) \models \mathrm{GCH} \}
\]
と定義する。

GCH は任意の無限基数 $\kappa$ に対し
\[
\mathcal{P}_\kappa(S) = \kappa^+
\]
を要求する。

\begin{theorem}[ZFC+GCH の Solver 制約]
特異基数 $\kappa$ に対し、$S \in \mathcal{Z}_{\mathrm{GCH}}$ であるためには、Solver $S$ は次の制約を同時に満たさなければならない：
\begin{enumerate}
  \item 任意の下位基数 $\kappa_i < \kappa$ に対する冪集合規則が一様であること。
  \item 極限段階における合成操作が $\kappa^+$ を超える新たな部分集合を生成しないこと。
  \item 異なる極限表現によって $\mathcal{P}_\kappa(S)$ の値が変化しないこと。
\end{enumerate}
このとき、これらの条件を満たす Solver は $\mathcal{Z}$ の中で極めて限られた部分集合を成す。
\end{theorem}

\subsection{2.7 定理の意味}
上記定理が示すのは、ZFC+GCH が特異基数において要求する冪集合の振る舞いが、Solver に対して非常に強い一様性条件を課すという事実である。これは ZFC 単体では現れない制約であり、GCH を追加することによって初めて発生する。
したがって、GCH は単なる基数算術の仮説ではなく、Solver の構成自由度を大幅に削減する選別条件として機能する。

\subsection{2.8 結論}
本稿では特異基数における $2^\kappa$ を Solver 集合上の写像として定式化し、その枠組みの中で ZFC+GCH が Solver 集合に極端な制約を課すことを定理として示した。
この結果は、特異基数の冪集合問題が未解決である理由を説明すると同時に、GCH が Solver 集合上で例外的な位置を占めることを明確にするものである。

\section*{3 補足事項：基数クラス $\mathrm{Card}$ について}
本稿で用いる $\mathrm{Card}$ は、ZFC における通常の基数（cardinal number）のクラスを表すものとする。
すなわち、$\mathrm{Card}$ は集合の濃度を同値類として表現した対象全体のクラスであり、
順序型としての基数、あるいは選択公理のもとで定義される標準的な基数概念を前提とする。

本稿において $\mathrm{Card}$ 自体の構造や性質を拡張・修正することは行わない。
$\mathrm{Card}$ はあくまで、Solver 集合上で定義される写像や評価の値域として用いられる補助的な対象である。

特に、$\mathcal{P}_\kappa(S) \in \mathrm{Card}$ のような表記は、
ある Solver $S$ のもとで許容される冪集合構成の結果として得られる濃度を、
ZFC 的意味での基数として記述しているに過ぎない。
これは、冪集合や基数が Solver に依存して再定義されることを意味するものではない。

\section{4 補足 : ZFCを解き直す}

\subsection{4.1 Ⅰ. 前提の再定義（ECA）}
\subsubsection{4.1.1 定義 1（Solver集合）}

Solver集合を次で定義する：
\[
\mathcal{S} = \{ S \mid S = (\mathcal{A}_S, \mathcal{R}_S) \}
\]

ここで $\mathcal{A}_S$を公理の選択集合、$\mathcal{R}_S$ を推論規則・構成制約の集合とする。
重要な点として、各Sは「公理そのもの」ではなく「ある公理の集まり」であること。

\subsubsection{4.1.2 定義 2（公理写像）}
各Solver集合 $S$ に対し、次のような写像を定義する：

\[
\mathcal{T} : \mathcal{S} \to \text{Th}
\]

ここで $\mathcal{T}(S)$ をSolver集合 $S$ によって表現される公理集合とし、矛盾の場合は未定義とする。

\subsubsection{4.1.3 定義 3（公理選択空間の測度）}

拡張連続性公理に基づき、次の測度式を定義する：

\[
I : \mathcal{S} \to [0,1]
\]

ここで $I(S)$ は $\mathcal{T}(S)$ が無矛盾であり、かつ集合演算に関して閉じて定義可能であることを表す。


\subsubsection{4.1.4 定義 4（極大 Solver集合）}

\[
S \in \mathcal{S}
\text{ が極大である} \iff
\forall S' \in \mathcal{S},\;
S \subsetneq S' \Rightarrow I(S') \le I(S)
\]

これは純粋に「公理の追加によって評価値が上がらない」という定義である。

\subsection{4.2 Ⅱ. 拡張連続性公理上におけるZFCの再定義}

\subsubsection{定義 5（ZFC_ECA）}
\[
\text{ZFC\_ECA} = \{ S \in \mathcal{S} \mid S \text{ は極大であり}\; \mathcal{T}(S) \text{ が ZFC の公理集合と同型} \}
\]

ここで重要なのは、まず「ZFCは前提ではない」ということ、次にZFCは「Solver集合 $S$ における測度 $\mathcal{I}$ に関する極大点として定義される」ということである。

\subsection{4.3 Ⅲ. 各公理の解釈}
\subsubsection{4.3.1 外延性公理}
外延性を含む Solver集合 $S$ は次を満たす：

\[
\forall x,y \in \mathcal{T}(S),\;
\left(\forall z\,(z \in x \leftrightarrow z \in y) \Rightarrow x = y \right)
\]

これは $\mathcal{T}(S)$ において「全く同じ要素からなる2つの集合は等しい」ことを主張する。
同値だが異なる要素を同時に保持するSolver集合は測度上において値が低下する。

\subsubsection{4.3.2 空集合公理}

空集合公理を含む Solver集合は次の式 : 

\[
\exists x\,\forall y\,(y \notin x)
\]

を満たす。

これは $\mathcal{T}(S)$ において、要素を持たない集合が存在することを要求する。
すなわち集合における最小元が存在することを表す。
Solver集合においては、空集合公理は「初期段階が $\mathcal{T}(S)$ 内で閉じて定義可能である」ということである。

\subsubsection{4.3.3 ペア集合公理}

Solver集合に次式を用いると : 

\[
\forall x,y\,\exists z\,\forall w
(w \in z \leftrightarrow (w=x \lor w=y))
\]

これは有限集合が $\mathcal{T}(S)$ 上で閉じている集合であることを要求する。

\subsubsection{4.3.4 和集合公理}
和集合公理で考える場合では次の式が成り立つ。

\[
\forall x\,\exists y\,\forall z
(z \in y \leftrightarrow \exists u (z \in u \land u \in x))
\]

これは $\mathcal{T}(S)$ において、集合 $x$ に属する各集合の要素全体を集めた集合が存在することを要求する。
すなわち、集合の集合に対して一次の和集合操作が閉じて定義可能であることを表す。
Solver集合においては、和集合公理は集合における階数の一段階の降下操作が $\mathcal{T}(S)$ 内で許容されることを意味する。

\subsubsection{4.3.5 冪集合公理}
冪集合公理においては次の式が成り立つ。

\[
\forall x\,\exists y\,\forall z (z \in y \leftrightarrow z \subseteq x)
\]

これはSolver集合に置き換えて考えた場合、部分集合全体が $\mathcal{T}(S)$ 集合上において対象となる集合が存在することを指す。

\subsubsection{4.3.6 無限公理}
無限集合では次式が成り立つ。

\[
\exists x(\varnothing \in x \land \forall y(y \in x \Rightarrow y \cup \{y\} \in x))
\]

これは有限閉包しか持たないSolver集合を排除する。


\subsubsection{4.3.7 分出・置換公理}
どちらも、定義可能な写像または述語に基づいて集合の存在が $\mathcal{T}(S)$ 内に制限されることを要求する。
無制限な集合の成立を防ぐため、原則として公理図式として定式化される。

\subsubsection{4.3.8 正則性公理}
また、正則性公理においては次の式が成り立つ。

\[
\forall x (x \neq \varnothing \Rightarrow \exists y (y \in x \land x \cap y = \varnothing))
\]

これは、集合の「無限降下列」（例: $A_{1}\ni A_{2}\ni A_{3}\ni \dots $）や、「自分自身を含む集合」（例: $A\in A$）といったパラドックスを引き起こす数学的矛盾を生み出さないことを集合に要求する。

\subsubsection{4.3.9 選択公理}
選択公理では次式が成り立つ。

\[
\forall X \bigl((\forall x \in X, x \neq \varnothing) \Rightarrow \exists f : X \to \bigcup X \bigr)
\]

これは $\mathcal{T}(S)$ において空でない集合族に対し、それぞれの要素を一つずつ選ぶ集合操作が定義可能であることを要求する。
Solver集合においては、選択公理は $\mathcal{T}(S)$ 内で定義可能な写像によって集合族の直積的構成が実現できることを意味する。

\end{document}